\documentclass[french]{book}
\usepackage{teach}
\usepackage{verbatim}
\usepackage{amsmath,amsfonts,amssymb}
\usepackage{pgfplots}
\pgfplotsset{compat=newest}
%\usepackage{geometry}
%\geometry{top=1cm, bottom=1cm, left=2cm, right=2cm}
\usepackage[utf8]{inputenc}
\usepackage[french]{babel}
\redefinecolor{thm}{title}{bgcolor}{alizarin}
\redefinecolor{prop}{title}{bgcolor}{meatbrown}
\redefinecolor{defn}{title}{bgcolor}{olivine}
\definecolor{alizarin}{rgb}{0.82, 0.1, 0.26}
\definecolor{meatbrown}{rgb}{0.9, 0.72, 0.23}
\definecolor{olivine}{rgb}{0.6, 0.73, 0.45}
\usepackage{mathrsfs}
\usetikzlibrary{arrows}


\begin{document}
\chapter{Géométrie et optimisation}

\section{Projeté orthogonal d'un point sur une droite}

\begin{defn}
Soit A un point du plan et $\Delta$ une droite.
Le \underline{projeté orthogonal} du point A sur la droite $\Delta$  est le point d'intersection de la droite $\Delta$ et de la perpendiculaire à $\Delta$ passant par A, noté H.
\end{defn}

\begin{prop}
Le projeté orthogonal du point A sur la droite $\Delta$ est le point de $\Delta$ le plus proche de A : cela signifie que, pour tout point M de $\Delta$ distinct de H, on a AM $>$ AH.
\end{prop}

Autrement dit, la distance entre A et un point de $\Delta$ est minimal lorsque ce point est H.

\begin{demo}
Soit M un point de $\Delta$ distinct de H. Comme AHM est un triangle rectangle, on a d'après le théorème de Pythagore.
$$AM^2=AH^2+HM^2$$ Puisque M est distinct de H on a nécessairement $HM^2>0$, donc en ajoutant $AH^2$ de part et d'autre de l'inégalité on obtient l'inégalité suivante,
$$AH^2+HM^2 > AH^2$$ mais d'après notre première égalité le premier membre de l'inégalité est égal à $AM^2$, d'où la dernière inégalité qui conclut la preuve,
$$AM^2>AH^2$$
ce qui est équivalent encore à 
$$AM>AH$$
\end{demo}


\begin{defn}
La longueur AH s'appelle la \underline{distance} du point A à la droite $\Delta$.
\vspace{4mm}
\end{defn}

\section{Trigonométrie}
\begin{center}
\definecolor{qqwuqq}{rgb}{0,0.39215686274509803,0}
\definecolor{xdxdff}{rgb}{0.49019607843137253,0.49019607843137253,1}
\definecolor{ududff}{rgb}{0.30196078431372547,0.30196078431372547,1}
\begin{tikzpicture}[scale=1.5,line cap=round,line join=round,>=triangle 45,x=1cm,y=1cm]
\clip(-1.2997560756662618,-0.8688016433538822) rectangle (7.487174934667167,3.9296864630703707);
\draw [shift={(-0.12228148239489928,-0.145589277575063)},line width=1pt,color=qqwuqq,fill=qqwuqq,fill opacity=0.10000000149011612] (0,0) -- (0.12455584620725263:0.20594369555468975) arc (0.12455584620725263:25.481334490023716:0.20594369555468975) -- cycle;
\draw[line width=1pt,color=qqwuqq,fill=qqwuqq,fill opacity=0.10000000149011612] (5.987723236966333,0.013317869151515865) -- (5.842099397398288,0.013001295587237477) -- (5.842415970962566,-0.13262254398080772) -- (5.988039810530611,-0.13230597041652933) -- cycle; 
\draw [line width=1pt] (-0.12228148239489928,-0.145589277575063)-- (5.988039810530611,-0.13230597041652933);
\draw [line width=1pt] (-0.12228148239489928,-0.145589277575063)-- (5.981744707757667,2.7634413051377784);
\draw [line width=1pt] (5.981744707757667,2.7634413051377784)-- (5.988039810530611,-0.13230597041652933);
\begin{scriptsize}
\draw [fill=ududff] (-0.12228148239489928,-0.145589277575063) circle (1.5pt);
\draw[color=ududff] (-0.08468827189359232,0.006459062753545894) node {$A$};
\draw [fill=ududff] (5.988039810530611,-0.13230597041652933) circle (1.5pt);
\draw[color=ududff] (6.145569065784339,0.013323852605368863) node {$B$};
\draw [fill=xdxdff] (5.981744707757667,2.7634413051377784) circle (1.5pt);
\draw[color=xdxdff] (6.038704275932517,2.9102651700746605) node {$C$};
\draw[color=qqwuqq] (0.3615230684749022,-0.02100009665374597) node {$\alpha$};
\end{scriptsize}
\end{tikzpicture}
\end{center}

\begin{thm}
Soit $ABC$ un triangle rectangle en $B$ et $\widehat{CAB}=\alpha$ alors on a $\cos(\alpha)^2+\sin(\alpha)^2=1$
\end{thm}

\begin{demo}
Soit $ABC$ un triangle rectangle en $B$ et $\widehat{CAB}=\alpha$  on sait que $\cos(\alpha)=\frac{AB}{AC}$ et $\sin(\alpha)=\frac{BC}{AC}$ donc $$\cos(\alpha)^2+\sin(\alpha)^2= \frac{AB^2 + BC^2}{AC^2}.$$ 
Or le théorème de Pythagore affirme que $AB^2 + BC^2=AC^2$, par suite $$\cos(\alpha)^2+\sin(\alpha)^2=\frac{AC^2}{AC^2}=1.$$
\end{demo}

\section{Problème d'optimisation géométrique}

Dans de nombreux problèmes, on cherche à rendre minimale ou maximale une grandeur géométrique (longueur, aire, volume...) : ce sont des problèmes d'\underline{optimisation} de grandeurs.


\begin{prop}
A, B et C sont trois points du plan.\\
Les trois propositions suivantes sont équivalentes: 
\begin{itemize}
\item C appartient au segment $[AB]$
\item La somme des distances AC+CB est minimale.
\item AB=AC+CB
\end{itemize}
\end{prop}




\end{document}
